\chapter{Appendix}

\section{Submission Artifact Inventory}

The final submission package is composed of three coordinated artifacts:

\begin{enumerate}
  \item \textbf{This final report PDF}, which documents requirements, architecture, project management,
  and proof-of-concept alignment.
  \item \textbf{The Git repository}, which contains the runnable MVP, root orchestration scripts,
  OpenAPI/contract notes, tests, Docker Compose configuration, CI workflow, and demo runbook.
  \item \textbf{The recorded presentation}, which should cover team responsibilities, key design
  decisions, and the proof-of-concept/demo flow.
\end{enumerate}

\section{Assignment Coverage Map}

Table~\ref{tab:coverage-map} maps the assignment brief's required deliverables to the final report and the
repository artifacts. This is included deliberately to reduce any ambiguity between the rubric, the
written report, and the delivered MVP.

\begin{longtable}{p{4.1cm}p{5.0cm}p{5.2cm}}
\caption{Assignment coverage map}\label{tab:coverage-map}\\
\toprule
\textbf{Assignment Requirement} & \textbf{Covered in Report} & \textbf{Covered in Repository / MVP} \\
\midrule
\endfirsthead
\toprule
\textbf{Assignment Requirement} & \textbf{Covered in Report} & \textbf{Covered in Repository / MVP} \\
\midrule
\endhead
Stakeholder analysis & Section 3, Stakeholder Analysis & Reflected in RBAC, AI policy, and deployment decisions \\
Functional requirements across collaboration, AI, document management, and user management & Section 3, Functional Requirements & Backend routes, frontend flows, realtime service, tests \\
Measurable non-functional requirements & Section 3, Non-Functional Requirements & Latency/error/reliability constraints reflected in design, scripts, and fallback behavior \\
At least ten user stories & Section 3, User Stories and Scenarios & Implemented or explicitly scoped in MVP alignment section \\
Requirements traceability matrix & Section 3, Traceability Matrix & Links stories to repo modules and services \\
Architectural drivers & Section 4, Architectural Drivers & Reflected in modular repo and service boundaries \\
C4 Level 1, 2, and 3 diagrams & Section 4, Figures and appendix source listings & Mermaid source in repo and final-report package \\
Feature decomposition & Section 4, Feature Decomposition & Monorepo service layout and domain modules \\
AI integration design & Section 4, AI Integration Design & AI service abstraction, async jobs, provider selection \\
API design with concrete contracts & Section 4, API Design & Backend OpenAPI, contract notes, automated tests \\
Authentication and authorization rationale & Section 4, Authentication and Authorization & Login route, bearer tokens, RBAC, realtime write restrictions \\
Communication model & Section 4, Communication Model & Backend sessions + Yjs realtime service + frontend reconnect behavior \\
Code structure and repository organization & Section 4, Code Structure \& Repository Organization & Root scripts, monorepo folders, shared code, tests \\
Data model and ER diagram & Section 4, Data Model & SQLite schema, repositories, versioning, AI jobs, export jobs \\
Four ADRs & Section 4, Architecture Decision Records & Reflected in code and tooling choices \\
Team structure and ownership & Section 5, Team Structure and Ownership & PR/branch history and service ownership in repo \\
Development workflow and methodology & Section 5, Development Workflow and Methodology & Branch structure, CI, root scripts, reviewable commit history \\
Risk assessment and milestones & Section 5, Risk Assessment and Timeline & Regression tests, demo reset, Docker Compose, CI, milestone table \\
PoC front-end/back-end communication & Section 6, Proof of Concept & Working MVP and automated smoke tests \\
README and runnable repository & Section 6, Repository and README Evidence & Root README, service READMEs, DEMO.md, Docker Compose \\
Editable diagram sources & Appendix Mermaid source listings & `docs/diagrams/*.mmd`, `final-report/figures/mermaid-src/*.mmd` \\
\bottomrule
\end{longtable}


\section{Version-Control and Collaboration Evidence}

The proof-of-concept rubric explicitly evaluates whether the Git history is clean and whether there is
evidence of team collaboration. The delivered repository supports that in three ways:

\begin{itemize}
  \item feature work was kept on named branches and integrated through pull-request-style review rather
  than as one unstructured dump into \texttt{main},
  \item commits were intentionally kept small and purposeful so architectural evolution is visible in the
  history,
  \item the repository contains both implementation artifacts and editable documentation artifacts so
  architecture changes can be versioned alongside code rather than maintained separately.
\end{itemize}

This matters for the assignment because the architecture, repo organization, and PoC are meant to reflect
the same design process rather than three unrelated deliverables.

\section{Demo Checklist}

The recommended live demo sequence is:

\begin{enumerate}
  \item Run \texttt{npm run demo:reset} from the repository root.
  \item Open the frontend at \texttt{http://localhost:5173}.
  \item Sign in as \texttt{user\_1} and create a new document.
  \item Open the same document in a second browser/tab and switch the user to \texttt{user\_2}.
  \item Grant editor access from the owner account.
  \item Type in both windows and show live sync and presence.
  \item Pause briefly to show autosave returning the header to a saved state.
  \item Trigger an AI action, wait for the async job, then demonstrate full apply or partial apply.
  \item Open version history and demonstrate revert.
  \item Export the document and show Swagger/OpenAPI at \texttt{http://localhost:3000/docs/}.
\end{enumerate}

\section{Mermaid Diagram Source Listings}

The assignment requires editable diagram source files. The report therefore includes Mermaid source
listings for the reconstructed architecture diagrams used as living artifacts.

\subsection{C4 Level 1 --- System Context (Mermaid Source)}
\lstinputlisting{figures/mermaid-src/c4-level-1-system-context.mmd}

\subsection{C4 Level 2 --- Container Diagram (Mermaid Source)}
\lstinputlisting{figures/mermaid-src/c4-level-2-container.mmd}

\subsection{C4 Level 3 --- Backend API Component Diagram (Mermaid Source)}
\lstinputlisting{figures/mermaid-src/c4-level-3-backend-components.mmd}

\subsection{Entity-Relationship Diagram (Mermaid Source)}
\lstinputlisting{figures/mermaid-src/er-data-model.mmd}

\section{Additional Mermaid Source Artifacts Available in the Repository}

The repository also includes supplementary Mermaid sources for dynamic interaction views:

\begin{itemize}
  \item \texttt{final-report/figures/mermaid-src/sequence-document-open-and-sync.mmd}
  \item \texttt{final-report/figures/mermaid-src/sequence-ai-job-flow.mmd}
\end{itemize}

These are maintained as editable assets but are intentionally not expanded into full code listings in the
main report PDF, because the assignment's core diagram requirement is already satisfied through the C4
and ER views included above.
